\section{总结与延伸}

Lecture 13 的核心是把“数据”从抽象名词还原成生产系统。网页不是直接可训练文本，Common Crawl 不是高质量数据，许可证不是一行元数据，过滤器不是中立工具。数据工程涉及来源、转换、过滤、去重、版权、隐私、质量和长期维护。

\begin{importantbox}{最终 takeaways}
\begin{enumerate}
    \item Data does not fall from the sky；数据获取和清洗需要大量工程和人类判断。
    \item 模型公司常公开架构和训练流程，却隐藏数据细节，说明数据是核心竞争力。
    \item Public web 受动态内容、认证、robots.txt、ToS、rate limit 和法律限制约束。
    \item Copyright/fair use/licensing 是训练数据管线不可回避的一部分。
    \item Common Crawl 是 raw source，不是成品训练集；HTML extraction 和 filtering 会塑造模型。
    \item 现代数据集在规则过滤、模型过滤、合成重写、dedup、license 审计之间折中。
    \item 下一步数据课程自然会进入更细的 processing、mixing 和 post-training 数据设计。
\end{enumerate}
\end{importantbox}

\begin{knowledgebox}{拓展阅读}
\begin{itemize}
    \item Common Crawl, C4, CCNet, RefinedWeb, FineWeb.
    \item The Pile, Dolma, DCLM, Nemotron-CC, The Stack, CommonPile.
    \item Copyright Act, Creative Commons, fair use, Terms of Service.
    \item LLaMA, GPT-3, Gopher, OLMo 2, Tulu data reports.
\end{itemize}
\end{knowledgebox}

\end{document}
